\section{Introduction}

Tempo estimation finds how fast a piece of music is, in beats per minute (BPM), and beat tracking finds the times at which a listener would tap along. Both are basic tasks in music information retrieval and a first step for downbeat tracking, chord recognition and structure analysis \cite{tutorial}. Our group (Default Project 2, DP2\_B) set three goals in the design specification: build one complete method per member from scratch, compare these methods with existing tools, and study three known weak spots: changing tempo, degraded audio and the bar level (downbeats).

Each member owns one method and one study. Irai Kumaran Sivanesan works on tempo estimation and tempo changes, Nalluraj Babu on dynamic programming beat tracking and degraded audio, and Usha Sundaresan Manivannan on neural network beat tracking and downbeats. \MeIntro{} All code, result files and the scripts that make every number, table and figure in this report are in our repository.\footnote{\url{https://github.com/Iraikumaran/beat-tracking}, branch \texttt{Dev}.}

\section{Related Work}

Most tempo and beat systems first compute an onset detection function \cite{bello} and then look for periodicities in it. Gouyon et al. \cite{gouyon} compared tempo induction algorithms, introduced the Accuracy1 and Accuracy2 measures, and showed that choosing half or double the correct tempo is the most common error. Percival and Tzanetakis \cite{percival} proposed a streamlined estimator based on autocorrelation followed by cross-correlation with pulse trains. Ellis \cite{ellis} framed beat tracking as dynamic programming over an onset envelope, the method behind the beat tracker of librosa \cite{librosa}. For music with changing tempo, Grosche and M\"uller \cite{plp} estimate a local pulse in short windows.

Neural networks now give the best results. B\"ock et al. \cite{downbeat} predict beat and downbeat activations with recurrent networks and decode them with a dynamic Bayesian network (DBN), as implemented in madmom \cite{madmom}. Krebs et al. \cite{krebsbar} find downbeats from beat-synchronous features. Davies and B\"ock \cite{tcn} showed that temporal convolutional networks (TCNs) are accurate and small, and Heydari et al. \cite{heydari} reduced decoding cost with a one dimensional state space. Davies et al. \cite{tutorial} give a practical overview of all of these methods. Mauch and Ewert \cite{degradation} showed that MIR systems, beat trackers included, lose accuracy under common audio degradations.

Evaluation follows Davies et al. \cite{daviseval} with the reference implementation in mir\_eval \cite{mireval}. GTZAN was introduced by Tzanetakis and Cook \cite{gtzan}; Sturm \cite{sturm} documents its faults, and Kereliuk et al. \cite{kereliuk} proposed the artist-filtered split we use. The beat and downbeat annotations come from Marchand et al. \cite{marchand}. GiantSteps Tempo \cite{giantsteps} is the second dataset in our plan.

\section{Data and Evaluation}

We use the GTZAN audio with the beat, downbeat and tempo annotations of the gtzan\_tempo\_beat repository \cite{tutorial, marchand}. All results are on the artist-filtered test split of Kereliuk et al. \cite{kereliuk} (\NTempo{} clips of 30\,s, all ten genres). The neural network is trained on 160 clips from the matching training and validation splits (16 per genre), so no test artist is seen in training. The studies on tempo changes and degraded audio use an evenly spaced subset of \NDeg{} test clips (every fifth clip, all genres). The reference tempo of a clip is $T = 60/\mathrm{median}(\Delta b)$, where $\Delta b$ are the annotated inter-beat intervals.

For a tempo estimate $\hat{T}$, Accuracy1 counts $|\hat{T}-T| \le 0.04\,T$, and Accuracy2 also accepts $\hat{T}$ within 4\% of $kT$ for $k \in \{2, 3, \tfrac{1}{2}, \tfrac{1}{3}\}$. For beats and downbeats we report the F-measure with a $\pm 70$\,ms tolerance window,
\begin{equation}
F = \frac{2PR}{P+R},
\end{equation}
where $P$ and $R$ are the precision and recall of the estimated beats, together with the continuity measures CMLt and AMLt \cite{daviseval}. CMLt asks for a continuous sequence at the correct metrical level, and AMLt also accepts double, half and off-beat sequences. Beats in the first 5\,s are ignored; for downbeats they are kept.

\textbf{Shared front end.} Audio is resampled to 22.05\,kHz and a short-time Fourier transform is computed with a 2048-sample Hann window and a hop of 256 samples ($f_r = 86.1$ frames per second). The magnitude is mapped to 40 mel bands and compressed as $M(f,n) = \log(1 + 100\,X_{\mathrm{mel}}(f,n))$. The onset envelope is the half-wave rectified spectral flux
\begin{equation}
o(n) = \sum_f \max\big(0, M(f,n) - M(f,n-1)\big),
\end{equation}
with a moving average removed and the result normalised. Our own methods use only NumPy, SciPy and scikit-learn; librosa \cite{librosa} and madmom \cite{madmom} are used only as reference systems.

\textbf{Repeatability.} Every experiment was run twice from a clean results folder on the same machine. Apart from run times, the two runs gave identical result files, and they also matched the files that each member had produced on their own laptop (two of 6090 rows in the tightness sweep differ, both at $\alpha = 3000$, by floating point effects).
